\textbf{R:} Dado que los atributos dados son variables continuas, nos apoyaremos de Bayes Naïve Gaussiano para determinar la clase a partir de las observaciones, por lo que usaremos la siguiente fórmula

Comenzaremos por obtener las probabilidades a priors de cada clase, entonces tenemos:

\begin{align*}
    P(Y=0)& = \frac{4}{9} = 0.444 \\
    P(Y=1) & = \frac{5}{9} = 0.556 \\
\end{align*}

% \mu es la media
Ahora procedemos a obtener las medias de cada clase para cada dato, partiendo de la fórmula $\frac{1}{N} \sum_{x} x$ entonces tenemos

\begin{align*}
    \mu_{y=0,x=1} & = \frac{0.5+0.3+0.2+0.4}{4} = 0.35 \\
    \mu_{y=0,x=2} & = \frac{0.6+0.6+0.5+0.6}{4} = 0.575 \\
    \mu_{y=0,x=3} & = \frac{0+0.01+0+0.1}{4} = 0.0275 \\\\
    \mu_{y=1,x=1} & = \frac{0.6+0.9+0.8+0.95+0.5}{5} = 0.75 \\
    \mu_{y=1,x=2} & = \frac{0.8+0.7+0.65+0.9+0.6}{5} = 0.73 \\
    \mu_{y=1,x=3} & = \frac{0.7+1+0.9+1+0.5}{5} = 0.82 \\
\end{align*}

% \sigma^2 es varianza de los datos
Ahora procedemos a obtener las varianzas de cada clase para cada dato, partiendo de la fórmula $\frac{1}{N-1} \sum_{x} (x-\mu_{y,x})^2$, entonces tenemos
\begin{align*}
    \sigma^2_{y=0,x=1} & = \frac{(0.5-0.35)^2 + (0.3-0.35)^2 + (0.2-0.35)^2 + (0.4-0.35)^2}{3} = 0.0167 \\
    \sigma^2_{y=0,x=2} & = \frac{(0.6-0.575)^2 + (0.6-0.575)^2 + (0.5-0.575)^2 + (0.6-0.575)^2}{3} = 0.0025 \\
    \sigma^2_{y=0,x=3} & = \frac{(0-0.0275)^2 + (0.01-0.0275)^2 + (0-0.0275)^2 + (0.1-0.0275)^2}{3} = 0.00235
\end{align*}
\begin{align*}
    \sigma^2_{y=1,x=1} & = \frac{(0.6-0.75)^2 + (0.9-0.75)^2 + (0.8-0.75)^2 + (0.95-0.75)^2 + (0.5-0.75)^2}{4} = 0.0375 \\
    \sigma^2_{y=1,x=2} & = \frac{(0.8-0.73)^2 + (0.7-0.73)^2 + (0.65-0.73)^2 + (0.9-0.73)^2 + (0.6-0.73)^2}{4} = 0.0145 \\
    \sigma^2_{y=1,x=3} & = \frac{(0.7-0.82)^2 + (1-0.82)^2 + (0.9-0.82)^2 + (1-0.82)^2 + (0.5-0.82)^2}{4} = 0.047 \\
\end{align*}

Ahora obtenremos el modelo de Bayes Naïve Gaussiano, para esto usaremos la fórmula \[P(x|y) = \frac{1}{\sqrt{2\pi\sigma^2_y}} e^{-\frac{(x-\mu_y)^2}{2\sigma^2_y}}\] 
Después aplicaremos la fórmula para la probabilidad conjunta
\[ P(y, x_1, x_2, ... , x_n) = P(y) \prod_{i=1}^{n} P(x_i | y) \]

\textbf{Nota}: Para calcular las distribuciones se utilizó programa en Python que se adjunta en la entrega de esta tarea.

\newpage
\begin{itemize}
    \item Ejemplo de evaluación 0:
    Comenzamos por obtener $P(x|y)$ de cada dato para la clase $y=0$
    \begin{align*}
        P(x_1=0.3|y=0) & = \frac{1}{\sqrt{2\pi(0.0167)}} e^{-\frac{(0.3-0.35)^2}{2(0.0167)}} = 2.87 \\
        P(x_2=0.7|y=0) & = \frac{1}{\sqrt{2\pi(0.0025)}} e^{-\frac{(0.7-0.575)^2}{2(0.0025)}} = 0.35 \\
        P(x_3=0.02|y=0) & = \frac{1}{\sqrt{2\pi(0.00235)}} e^{-\frac{(0.02-0.0275)^2}{2(0.00235)}} = 8.11
    \end{align*}

    Calculamos la probabilidad conjunta
    \begin{align*}
        P(y=0, x_1=0.3, x_2=0.7, x_3=0.02) & = P(y=0) \cdot P(x_1=0.3|y=0) \cdot P(x_2=0.7|y=0) \cdot P(x_3=0.02|y=0) \\
        & = 0.444 \cdot 2.87 \cdot 0.35 \cdot 8.11 \\
        & = 3.61
    \end{align*}

    Ahora veamos que sucede en la clase $y=1$
    \begin{align*}
        P(x_1=0.3|y=1) & = \frac{1}{\sqrt{2\pi(0.0375)}} e^{-\frac{(0.3-0.75)^2}{2(0.0375)}} = 0.138 \\
        P(x_2=0.7|y=1) & = \frac{1}{\sqrt{2\pi(0.0145)}} e^{-\frac{(0.7-0.73)^2}{2(0.0145)}} = 3.211 \\
        P(x_3=0.02|y=1) & = \frac{1}{\sqrt{2\pi(0.047)}} e^{-\frac{(0.02-0.82)^2}{2(0.047)}} = 0.002
    \end{align*}

    Calculamos la probabilidad conjunta
    \begin{align*}
        P(y=1, x_1=0.3, x_2=0.7, x_3=0.02) & = P(y=1) \cdot P(x_1=0.3|y=1) \cdot P(x_2=0.7|y=1) \cdot P(x_3=0.02|y=1) \\
        & = 0.556 \cdot 0.138 \cdot 3.211 \cdot 0.002 \\
        & = 0.00049 
    \end{align*}
    
    Veamos que la probabilidad conjunta de la clase $y=0$ es mayor que la de la clase $y=1$, por lo que se estima que el ejemplo de evaluación 0 pertenece a la clase $y=0$ con una probabilidad de $\frac{3.61}{3.61 + 0.00049} = 0.9998$.

    \newpage
    \item Ejemplo de evaluación 1:
    \newpage
    \item Ejemplo de evaluación 2:
    \newpage
    \item Ejemplo de evaluación 3:
    \newpage
    \item Ejemplo de evaluación 4:
\end{itemize}